\documentclass{report}

\begin{document}
  En semestres anteriores hemos aprendido un poquito acerca de analisis matematico, esto es en materias de calculo
  pues lo que aprendemos entonces es una estudio centrado en $R$ y $R^2$. Sin embargo, la intencion del analisis matematico
  es extender estos conceptos a $R^n$ y esto conlleva a generalizar y abstraer muchos conceptos que cuya interpretacion
  geometrica y analitica es mas sencilla de procesar puesto que podemos ver "graficamente" que esta ocurriendo. Un ejemplo de ello
  es la nocion de distancia, y para entenderlo mejor debemos preguntarnos
  
  ¿Que es una distancia?
  
  ¿Que caracteriza a una distancia?

  Ahora bien, existen dos metricas que podemos remarcar, una es la clasica metrica sobre $R$; el valor absoluto definido como:
  $$d(x,y) = \mid x-y \mid \; \forall \; x,y \in R$$

  Por otro lado, una metrica muy usada y caracteristica que mas adelante se va a estudiar es la metrica discreta, esta esta definida por: 
  
  $d(x,y)^{2y} = 1$ si $x \neq y$ 
  
  $d(x,y) = 0$ si $x = y$ 
  

  Es importante notar que esta metrica se puede definir claramente por $d: R \times R \to {0,1}$
\end{document}

